\section{Evaluasi}
\label{sec:evaluasi-bab4}

\subsection{Soal Pilihan Ganda}

\begin{enumerate}
    \item Unit yang bertanggung jawab menghitung alamat fisik 20-bit
          dari pasangan segmen:offset pada 8086 adalah \dots
          \begin{enumerate}
              \item \eng{Execution Unit}
              \item \eng{Bus Interface Unit}
              \item \eng{Arithmetic Logic Unit}
              \item register \eng{flag}
          \end{enumerate}

    \item Ukuran \eng{instruction queue} pada mikroprosesor 8086
          adalah \dots
          \begin{enumerate}
              \item 2 byte \item 4 byte
              \item 6 byte \item 16 byte
          \end{enumerate}

    \item Bila \reg{CS} $=$ \hex{1000} dan \reg{IP} $=$ \hex{0025},
          maka instruksi berikutnya diambil dari alamat fisik \dots
          \begin{enumerate}
              \item \hex{10025} \item \hex{10250}
              \item \hex{1002} \item \hex{100025}
          \end{enumerate}

    \item Sinyal yang menandai bahwa pin \reg{AD0}--\reg{AD15}
          sedang membawa alamat sehingga nilainya perlu ditangkap
          \eng{latch} adalah \dots
          \begin{enumerate}
              \item \reg{READY} \item \reg{ALE}
              \item \code{M/IO} \item \reg{RESET}
          \end{enumerate}

    \item Setelah \reg{RESET} aktif, 8086 memulai eksekusi dari
          alamat fisik \dots
          \begin{enumerate}
              \item \hex{00000} \item \hex{10000}
              \item \hex{FFFF0} \item \hex{FFFFF}
          \end{enumerate}
\end{enumerate}

\subsection{Soal Uraian}

\begin{enumerate}
    \item Gambarkan diagram blok 8086 yang memuat EU, BIU, bus
          internal, dan bus eksternal. Sebutkan register yang berada
          di dalam masing-masing unit dan jelaskan mengapa EU tidak
          boleh mengakses bus luar secara langsung.

    \item Jelaskan mekanisme \eng{pipelining} pada 8086 dengan
          bantuan contoh tiga instruksi. Hitung perbedaan jumlah
          siklus antara eksekusi tanpa antrean dan dengan antrean,
          lalu jelaskan kapan keuntungan tersebut hilang.

    \item Sebuah instruksi \lstasm|MOV AX, [BX]| memerlukan waktu
          lebih lama daripada \lstasm|MOV AX, BX|. Uraikan tahapan
          siklus instruksi untuk kedua instruksi tersebut dan
          tunjukkan pada tahap mana perbedaan waktunya muncul.

    \item Diketahui tiga instruksi berturut-turut bertempat pada
          alamat \hex{10000}, \hex{10003}, dan \hex{10005} dengan
          panjang berturut-turut 3 byte, 2 byte, dan 2 byte. Hitung
          nilai \reg{IP} setelah ketiga instruksi selesai bila
          \reg{CS} $=$ \hex{1000} dan \reg{IP} mula-mula \hex{0000}.
          Jelaskan hubungan antara panjang instruksi dan kenaikan
          \reg{IP}.
\end{enumerate}

\subsection{Kisi-kisi Penilaian}

\begin{table}[H]
    \centering
    \caption{Kisi-kisi evaluasi Bab~\ref{chap:arsitektur-8086}.}
    \label{tab:kisi-bab4}
    \small
    \begin{tabular}{@{}p{1.5cm}p{5.0cm}p{2.6cm}c@{}}
        \toprule
        \textbf{Sub-CPMK} & \textbf{Indikator yang diukur} &
        \textbf{Bentuk soal} & \textbf{Bobot} \\
        \midrule
        1.1 & Menjelaskan pembagian tugas BIU dan EU beserta
        register masing-masing & Pilihan ganda 1; Uraian 1 & 30\% \\
        \addlinespace
        1.1 & Menjelaskan \eng{instruction queue} dan
        \eng{pipelining} instruksi & Pilihan ganda 2; Uraian 2 & 25\% \\
        \addlinespace
        1.1 & Menjelaskan bus alamat, data, kontrol dan sinyal
        kendali & Pilihan ganda 4, 5 & 20\% \\
        \addlinespace
        1.1 & Menghitung alamat fisik dan menelusuri siklus
        instruksi & Pilihan ganda 3; Uraian 3, 4 & 25\% \\
        \bottomrule
    \end{tabular}
\end{table}

\begin{checklistbox}
\textbf{Checklist Pencapaian Kompetensi Bab~\ref{chap:arsitektur-8086}}

Seorang mahasiswa dinyatakan kompeten pada Sub-CPMK 1.1 apabila:
\begin{itemize}
    \item[\checked] dapat membedakan arsitektur internal dan
          eksternal serta menyebutkan contoh keduanya;
    \item[\checked] dapat menjelaskan pembagian tugas BIU dan EU
          beserta alasan pemisahan tersebut;
    \item[\checked] dapat menghitung alamat fisik dari pasangan
          segmen dan offset dengan benar;
    \item[\checked] dapat menjelaskan peran \eng{instruction queue}
          enam byte dan pengaruh \eng{pipelining} pada kinerja;
    \item[\checked] dapat menelusuri satu instruksi melalui tahapan
          \eng{fetch}--\eng{decode}--\eng{execute}--\eng{store}
          beserta perubahan \reg{CS}, \reg{IP}, dan antrean;
    \item[\checked] dapat menyebutkan fungsi sinyal \reg{ALE},
          \code{M/IO}, \reg{RD}, \reg{WR}, \reg{READY},
          \reg{RESET}, dan \code{MN/MX}.
\end{itemize}
\end{checklistbox}
